\editionPage[\pdfbookmark[0]{Corrigés des bilans}{bm-bilan-A-cor-auto-parent}]{bilan-A-cor-auto}{Bilan A · Corrigé}{Automatismes expliqués}{\hyperref[bilan-A-auto]{Énoncé p.\,\pageref*{bilan-A-auto}}}
\begin{blocEdition}{bilan-A-cor-auto}
\begin{questions}
\item Le produit est prioritaire : $7+12=\res{19}$.
\item $84:7=\res{12}$.
\item $74=9\times8+2$ : quotient $\res8$, reste $\res2$.
\item $\res{-3{,}4<-3{,}04}$ ; le premier est plus éloigné de zéro du côté négatif.
\item $-7+12=\res5$.
\item $8/12+3/12=\res{11/12}$.
\item $11\times11=\res{121}$.
\item $\res{4x+8}$ par distributivité.
\item $x=18/5=\res{3{,}6}$.
\item Un quart de 60 vaut $\res{15}$.
\item $180-55-65=\res{60^\circ}$.
\item Trois faces sur six : $\res{1/2}$.\end{questions}
\end{blocEdition}
\editionPage{bilan-A-cor-1}{Bilan A · Corrigé 1}{Une sortie et son budget}{\hyperref[bilan-A-ex-1]{Énoncé p.\,\pageref*{bilan-A-ex-1}}}
\begin{blocEdition}{bilan-A-cor-1}
\editionEtape{Dépenses}Les entrées coûtent $28\times4{,}50=126$ euros. Le total vaut $98+126=\res{224}$ euros.\editionEtape{Reste}Les sommes reçues valent $28\times8{,}50=238$ euros. Il reste $\res{14}$ euros, soit $14/238=1/17\approx\res{5{,}88\,\%}$ du total reçu.\editionEtape{Partage}Chaque atelier reçoit $14:7=\res2$ euros. Le pourcentage prend pour référence les sommes reçues, pas les dépenses.
\end{blocEdition}
\editionPage{bilan-A-cor-2}{Bilan A · Corrigé 2}{Un triangle et sa médiane}{\hyperref[bilan-A-ex-2]{Énoncé p.\,\pageref*{bilan-A-ex-2}}}
\begin{blocEdition}{bilan-A-cor-2}
\editionEtape{Aire entière}$10\times6:2=\res{30\ \mathrm{cm}^2}$.\editionEtape{Première médiane}CM partage le triangle en deux triangles de bases égales et de même hauteur : ACM a une aire de $\res{15\ \mathrm{cm}^2}$.\editionEtape{Second partage}Dans ACM, N est le milieu de [AM]. CN partage donc cette aire en deux : ACN et NCM ont chacun $\res{7{,}5\ \mathrm{cm}^2}$. On n’a pas besoin de connaître la longueur des médianes.\label{fig-bilan-A-cor-2}\figureGeo[corrige]{bilan-triangle}

\end{blocEdition}
\editionPage{bilan-A-cor-3}{Bilan A · Corrigé 3}{Une enquête et un tirage}{\hyperref[bilan-A-ex-3]{Énoncé p.\,\pageref*{bilan-A-ex-3}}}
\begin{blocEdition}{bilan-A-cor-3}
\editionEtape{Fréquences}Elles valent $\res{40\,\%,30\,\%,30\,\%}$.\editionEtape{Diagramme}Les angles valent $\res{144^\circ,108^\circ,108^\circ}$ ; leur somme est 360 degrés.\editionEtape{Tirage}Les fiches sont supposées mélangées et équiprobables. Six fiches sur vingt correspondent au bus : $\res{3/10}$. Cela décrit un tirage dans les fiches recueillies, pas nécessairement les habitudes de tout le collège.\label{fig-bilan-A-cor-3}\figureGeo[corrige]{bilan-enquete}

\end{blocEdition}
\editionPage{bilan-A-cor-4}{Bilan A · Corrigé 4}{Un bac et des récipients}{\hyperref[bilan-A-ex-4]{Énoncé p.\,\pageref*{bilan-A-ex-4}}}
\begin{blocEdition}{bilan-A-cor-4}
\editionEtape{Capacité totale}$40\times25\times30=30000$ cm$^3$, soit $\res{30\ \mathrm L}$.\editionEtape{Niveau choisi}Les $3/5$ de 30 L valent $(30:5)\times3=\res{18\ \mathrm L}$.\editionEtape{Bouteilles}$18:1{,}5=\res{12}$ bouteilles. Le quotient est entier : aucun reste de bouteille n’est nécessaire.\label{fig-bilan-A-cor-4}\figureGeo[corrige]{bilan-bac}

\end{blocEdition}
\editionPage{bilan-B-cor-auto}{Bilan B · Corrigé}{Automatismes expliqués}{\hyperref[bilan-B-auto]{Énoncé p.\,\pageref*{bilan-B-auto}}}
\begin{blocEdition}{bilan-B-cor-auto}
\begin{questions}
\item De gauche à droite : $8\times2=\res{16}$.
\item La somme des chiffres vaut 9 : $\res{\text{oui}}$.
\item Les deux valent $\res{2{,}5}$ ; la distance à zéro est la valeur absolue.
\item $5-6=\res{-1}$.
\item $3/4=6/8>5/8$.
\item $10/12-3/12=\res{7/12}$.
\item $27+4=\res{31}$.
\item $\res{5x+6}$.
\item $x=3-7=\res{-4}$.
\item $4500:1000=\res{4{,}5\ \mathrm L}$.
\item $27:3=\res9$.
\item $\res{\text{Oui}}$, comme les longueurs et les angles.\end{questions}
\end{blocEdition}
\editionPage{bilan-B-cor-1}{Bilan B · Corrigé 1}{Deux tarifs de location}{\hyperref[bilan-B-ex-1]{Énoncé p.\,\pageref*{bilan-B-ex-1}}}
\begin{blocEdition}{bilan-B-cor-1}
\editionEtape{Tableau}A donne $0,4,8,12,16$ ; B donne $6,8,10,12,14$ euros.\editionEtape{Comparer}Les tarifs sont égaux à $\res3$ h pour 12 euros. A est moins cher à 2 h ; B à 4 h.\editionEtape{Proportionnalité}Seul $\res A$ est proportionnel : coefficient 4 et droite passant par l’origine. B a un supplément fixe ; son prix à 0 h vaut 6 euros.\label{fig-bilan-B-cor-1}\figureGeo[corrige]{bilan-tarifs}

\end{blocEdition}
\editionPage{bilan-B-cor-2}{Bilan B · Corrigé 2}{Un parallélogramme et ses diagonales}{\hyperref[bilan-B-ex-2]{Énoncé p.\,\pageref*{bilan-B-ex-2}}}
\begin{blocEdition}{bilan-B-cor-2}
\editionEtape{Milieu commun}Les diagonales ont le même milieu : ABCD est un parallélogramme.\editionEtape{Nature plus précise}Des diagonales égales en font un rectangle. Si elles sont aussi perpendiculaires, il est en plus losange : c’est un carré.\editionEtape{Mesures du rectangle}Avec les seules dimensions 8 cm et 5 cm, $A=8\times5=\res{40\ \mathrm{cm}^2}$ et $P=2(8+5)=\res{26\ \mathrm{cm}}$. Ce rectangle précis n’a pas de diagonales perpendiculaires, car ses côtés sont différents : les cas de la question précédente étaient conditionnels.\label{fig-bilan-B-cor-2}\figureGeo[corrige]{bilan-rectangle}

\end{blocEdition}
\editionPage{bilan-B-cor-3}{Bilan B · Corrigé 3}{Un cylindre à remplir}{\hyperref[bilan-B-ex-3]{Énoncé p.\,\pageref*{bilan-B-ex-3}}}
\begin{blocEdition}{bilan-B-cor-3}
\editionEtape{Volume}$V=\pi\times3^2\times12=\res{108\pi\ \mathrm{cm}^3}$.\editionEtape{Capacité}Un cm$^3$ vaut un mL, donc environ $\res{339\ \mathrm{mL}}$.\editionEtape{Comparaison}La base est la même, les volumes sont proportionnels aux hauteurs. Le rapport vaut $8/12=2/3$, soit environ $\res{66{,}7\,\%}$. Utiliser les valeurs exactes évite les écarts dus aux arrondis.\label{fig-bilan-B-cor-3}\figureGeo[corrige]{bilan-cylindre}

\end{blocEdition}
\editionPage{bilan-B-cor-4}{Bilan B · Corrigé 4}{Un script à lire et à réparer}{\hyperref[bilan-B-ex-4]{Énoncé p.\,\pageref*{bilan-B-ex-4}}}
\begin{blocEdition}{bilan-B-cor-4}
\editionEtape{Lecture}Trois côtés d’un carré sont tracés ; le chemin n’est pas fermé.\editionEtape{Réparation}Mettre quatre répétitions. Le périmètre est $4\times40=\res{160}$ pas.\editionEtape{Autre taille}Pour un périmètre de 240, chaque côté doit mesurer $240:4=\res{60}$ pas. On garde quatre répétitions et des virages de 90 degrés. Des virages de 60 degrés ne donnent pas les angles droits du carré.\label{fig-bilan-B-cor-4}\figureGeo[corrige]{bilan-script}

\end{blocEdition}
